\section{Two horizons and the scalar spectral comparison}
\label{sec:dispersion}

The comparison needed for the trace polarization is stronger than the
order-two comparison furnished by its physical differential field.
We prove that the quadratic moment of the bounded-profile generators
satisfies an order-zero comparison on the vacuum wave $l$.
No local field with vacuum wave $l$ is used in this Section.

The argument has two parts.  First, the same causal stress-tensor
commutator is evaluated from two null directions.  This relates their
bounded-profile matrices after all local subtraction terms have been
removed.  Second, the difference between Hankel orders two and zero is
shown to have only one possible tail.  A logarithmically long radial
test makes the two-horizon identity detect the squared norm of that
tail, and positivity forces it to vanish.

We recall precisely the inputs from Sections~\ref{sec:foundations}
and~\ref{sec:null-current}.  On the internal Hilbert space at fixed
light-front coordinates, $h\geq0$ is the parabolic operator and
$h_\nu=-\partial_M^2-M^{-1}\partial_M+\nu^2M^{-2}$ has its regular
Hankel realization in $L^2(\R_+,M\,\dd M)$.  The comparison
\eqref{F:spectral-equation} applies with
\begin{equation}\label{D:orders}
\begin{array}{c|c}
 \text{constant polarization label}&\nu\\ \hline
 c_V=l+j_z,\quad c_T=l/3+s_z&0\\
 l,\ j_z,\ s_z&2\\
 j_i,\ s_{zi}&1.
\end{array}
\end{equation}
The bounded Ward formula \eqref{N:bounded-Ward}, including its polynomial
extension, and the moment representation give
\begin{equation}\label{D:measure-input}
 H_g=b\int g(Y+z')\,\mu(\dd z'),\qquad
 \mu(1)=\id,\qquad \mu(z'_i)=0,\qquad
 \kappa:=\mu(|z'|^2)\leq h.
\end{equation}
Here $\mu$ is a positive commutative operator-valued measure; its values
commute with $h$, $Y$ and $b$.  It need not be projection-valued.
The construction is covariant under dilatations and rotations, independent
of the boost coordinate, and agrees with the reflected construction.
Consequently, on every copy in \eqref{D:orders},
\begin{equation}\label{D:inherited-comparison}
 \mathbf1_{(s,\infty)}(\kappa)\mathbf1_{[0,s]}(h_\nu)=0,
 \qquad \norm{\kappa^t\psi}\leq\norm{h_\nu^t\psi}
 \quad(t>0).
\end{equation}
Indeed $\kappa$ and $h$ commute, so $\kappa\leq h$ combines with the
spectral comparison for $h$.  All internal formulas below are Hilbert
space formulas or distributions tested on compact output mass intervals;
they do not presume integral kernels for $\mu$.

In particular, the comparison in \eqref{D:inherited-comparison}
retains the embedding of each scalar copy in the ambient internal
space.  It bounds the spectrum of its image; it does not say that
$\kappa$ preserves that copy.  Reduction will be proved only after
the charge identities are available.

\subsection{The plane matrix and its causal continuation}

The bounded Ward identity initially supplies a localized plane matrix.
To compare the two horizons we need its causal continuation and a
fixed-normal restriction in the particular three-point channel used
here.  We establish these two operations before carrying out any
Cauchy subtraction.

Fix future null vectors $n,\bar n$ with $n\cdot\bar n=1$ and use
the transverse frame of Section~\ref{sec:foundations}.  Write
\begin{equation}\label{D:kinematics}
\begin{aligned}
 p&=a n+b\bar n+w, &p^2&=M_o^2>0,\\
 k&=a_x n+q, &p-k&=(a-a_x)n+b\bar n+r,\qquad w=q+r,\\
 X&=-k^2=|q|^2, &z&=-(p-k)^2=|r|^2-2b(a-a_x).
\end{aligned}
\end{equation}
The comparison point is $a_x=a$, where both virtualities are strictly
spacelike: $X=|q|^2>0$ and $z=|r|^2>0$.

\Needspace{6\baselineskip}
\begin{lemma}[The full-plane matrix]\label{D:plane-matrix}
On the spacelike-transfer part of $n\cdot k=0$, the full $T_{nn}$
matrix integrated along $n$ is the matrix of
$b\int g(Y+z')\mu(\dd z')$.  Initially the assertion is a joint
distributional identity, averaged over the null direction $n$ and tested
in the normal displacement $u$ of the plane. For the stress-tensor
three-point matrix with timelike spin-one output used below, the angular
average can be removed: $n$ may be fixed. The displacement $u$ remains
a distributional variable, whose Fourier transform supplies $a_x$.
\end{lemma}

\begin{proof}
Translate the plane normally by $u$, test $u$ with a compact smooth
function, and choose $\widehat g$ smooth with compact support in a
transverse annulus.  Average $n$ over a small smooth patch, taking a
smooth transverse frame and a sphere normalization of $n$.  Opposite
wedge tests can be chosen uniformly for these small changes.  Insert
a cutoff in the longitudinal coordinate $v$ and split its matrix by
the separating step used in \eqref{N:bounded-Ward}.  One term is that
Ward formula; the remaining term is an all-plane vacuum matrix.
The latter vanishes at spacelike transfer by the vacuum spectrum.
The operator on the Ward side at normal coordinate $u$ is
$U_a(u)H_gU_a(-u)$.

The removal of the $v$ cutoff has an independent distributional
meaning.  In momentum space it inserts $2\pi\delta(n\cdot k)$ after
averaging over $n$. On its support, $k=a_xn+q$ with $q$ in the
chosen annulus.  The derivative of $n\cdot k$ in an angular direction
of $n$ has the nonzero transverse component $q$.  The coarea formula
therefore turns the averaged delta insertion into a smooth multiplier.
The Fourier transform of the normal test is Schwartz.  Derivatives of
the moving coordinates cost polynomial factors, absorbed by this
decay and the compact annular test.  Thus the cutoff multipliers
converge in the Schwartz seminorms needed to pair with any of the
tempered external field distributions.  For an arbitrary fixed list
of original fields, these are the finite sums of homogeneous-list
distributions established in Section~\ref{sec:foundations}; no single
scaling dimension is assigned to the whole list.

The identity on opposite wedge tests extends to arbitrary external
translations by positive-energy tube uniqueness.  On the operator
side, $|H_g|\leq\norm{g}_\infty b$ supplies the polynomial energy
bound: insert positive-energy exponential damping in the two
external vectors, obtain the holomorphic matrix elements, and remove
the damping as distributional boundary values.  Equality on the open
wedge patch then determines the boundary distribution everywhere.
This establishes the identity jointly with the normal parameters.
The fixed-normal pullback needed here is justified in
Lemma~\ref{D:causal-continuation}.
\end{proof}

\Needspace{6\baselineskip}
\begin{lemma}[Fixed-spacelike-virtuality dispersion]
\label{D:causal-continuation}
Test $[T_{\alpha\beta}(x),T_{\rho\sigma}(0)]\Omega$ against the
spin-one output wave, with output measure $\dd^4p$ on the open
timelike cone.  At fixed $X>0$, the contracted matrix in
\eqref{D:kinematics} is the jump of polynomially bounded functions
holomorphic in the upper and lower $z$ half-planes.  The jump is the
first Wightman order and is supported in $z\leq0$.
At fixed $z>0$, the corresponding $X$ jump is the reversed Wightman
order with the same sign.  The two gap values come from one causal
splitting of this commutator.
\end{lemma}

\begin{proof}
Retain compact timelike output tests throughout the construction.
Testing against a $j$ label amounts to testing with $V$ and smooth
transverse projectors on that compact set. Lorentz and scale covariance
trivialize the open timelike orbit, leaving only the finite rest-rotation
representation. They identify the resulting distributions in simultaneous
moving coordinates. Fixed-$p$ formulas below refer to this covariant
disintegration; their estimates are uniform on compact positive output
mass ranges.

Local commutativity supports the relative-position commutator in
$\overline V_+\cup\overline V_-$.  Off the origin choose future and
past cutoffs constant near these two cones on the unit sphere.  Finite
distribution order allows their products to be extended through the
origin by subtracting a sufficiently long Taylor jet of the test
function. This is the cone-supported distribution extension used in
causal splitting; see \cite[Section~5.1]{EpsteinGlaser1973}.
Restore the total commutator by a distribution supported
at the origin.  This gives tempered cone-supported distributions
whose difference is the commutator.  Average the splitting over rest
rotations and transport it along the timelike Lorentz and scale
orbit.  Its local ambiguity is a finite sum of derivatives of a
delta function.  With the Fourier convention
$D(p,x)=\int e^{ikx}D(p,k)\,\dd k$, the resulting boundary values
$D^+,D^-$ are holomorphic in the upper and lower primitive tubes,
respectively, with polynomial bounds in momentum and inverse
distance to the tube boundary
\cite[Section~1, pp.~155--156]{BrosIagolnitzer1973}.

Here is a direct path from those tubes to the required half-planes.
In the rest frame of $p$, set
\begin{equation}\label{D:virtuality-path}
 k^0=\frac{M_o^2-X+z}{2M_o},\qquad
 R^2=(k^0)^2+X,\qquad \boldsymbol k=R\omega,
 \quad\omega\in S^2\text{ real}.
\end{equation}
Write $k^0=u+iv$, $R=s+it$.  Direct calculation gives
\begin{equation}\label{D:tube-distance}
 v^2-t^2=
 \frac{2Xv^2}{u^2+v^2+X+|(u+iv)^2+X|}>0,
 \qquad
 |v|-|t|\geq\frac{X|v|}{2(u^2+v^2+X)}.
\end{equation}
Thus \eqref{D:virtuality-path} lies strictly in the relevant tube
whenever $\operatorname{Im}z\ne0$.  On compact positive $X,M_o$
ranges the distance is bounded below by a constant times
$|\operatorname{Im}z|/(1+|z|)^2$.  The tube bounds therefore give
polynomial bounds in $|z|+|\operatorname{Im}z|^{-1}$.

The square root in \eqref{D:virtuality-path} introduces no extra
branch.  A tensor with the fixed finite collection of external
indices has only finitely many rotation harmonics.  A harmonic of
degree $s$ is $R^s$ times a Cartesian harmonic polynomial in
$\omega$; its remaining scalar coefficient is even and analytic in
$R$.  At $R=0$, where the path is still inside a tube, analyticity
in the Cartesian momentum components proves regularity of that
coefficient.  Divide by $R^s$ and reconstruct with Cartesian harmonic
polynomials.  The scalar coefficients are single-valued functions of
the two invariants.  They can consequently be evaluated at the
actual moving light-front momentum in \eqref{D:kinematics}, even
though that complex momentum itself has null imaginary part.
The branch points $k^0=\pm i\sqrt X$ are nonreal; none belongs
to the real locus on which a distributional restriction is taken.

We also verify that the real boundary obtained this way is the
required restriction.  For fixed real $\omega$,
$\dd R/\dd k^0=k^0/\sqrt{(k^0)^2+X}$ has absolute value less than
one.  The real tangent is strictly timelike.  Integration by parts
against a compact test interval therefore gives the uniform decay
on the dual cone needed for the Fourier--Laplace boundary limit.
The harmonic reconstruction gives the moving-direction restriction.
There is a useful equivalent noncharacteristic check.  The tube
boundary distributions have ordinary wavefront covectors in the
corresponding causal cones: compact localization in momentum
convolves the cone-supported tempered inverse Fourier transform
with a Schwartz function, giving rapid decay in directions outside
that cone.  For the full finite tensor family, rest-rotation
covariance gives equations $L_{ij}D=\rho_{ij}D$, where $L_{ij}$
is the angular rotation vector field and $\rho_{ij}$ a constant
matrix on the tensor indices.  The principal symbols of these
equations force every singular covector to annihilate the angular
orbit.  Since real spacelike $k$ has nonzero spatial part, its
annihilator is exactly $\Span\{p,k\}$.  If
$\xi=A p+B k$ annihilates the tangent $n$ in
\eqref{D:kinematics}, then $A b=0$, so $\xi=B k$.  As $k^2=-X<0$,
no nonzero such covector is causal. The pullback along
$a_x\mapsto a_xn+q$ is therefore noncharacteristic. On compact
positive $X$ and timelike output ranges the separation of its conormal
bundle from these wavefront cones persists in a small angular patch
of $n$. The restricted distribution thus depends smoothly on that
patch, removing the angular average in Lemma~\ref{D:plane-matrix}.
The normal displacement $u$ is still tested, or Fourier transformed;
no point evaluation of an arbitrary distribution in $u$ is needed.
This is a restriction
statement for this covariant three-point matrix, not for arbitrary
null-plane field distributions.

In the first Wightman order the vacuum leg has momentum $p-k$;
in the reverse order it has momentum $k$.  At $X>0$ the reverse
order vanishes by the spectrum condition, and the first can occur
only for $z\leq0$.  The two holomorphic boundary values glue in
the positive $z$ gap.  If instead $z>0$ is fixed and $X$ varies,
upper-half-plane continuation corresponds to the lower original
tube; its jump is minus the commutator, hence the reverse Wightman
order with a positive sign.  Both gap values are values of the
same split commutator at two spacelike virtualities.  This proves
the comparison directly from causal tubes, without a crossing
assumption.
\end{proof}

\subsection{Subtractions and the two-horizon identity}

Causal continuation identifies the two spacelike gap values, but a
dispersion formula determines each only up to a polynomial.  The next
step tracks the dependence of these polynomials on the transverse
variables.  The useful conclusion is their separated form: each
ambiguity is polynomial in at least one of the two variables.

Write $[\psi]_{j_A}$ for the output distribution tested against the
constant spin-one label in the indicated horizon frame, and put
$\mu_n(q)=\int e^{iq\cdot z'}\mu(\dd z')$ in that frame.
Application of Lemma~\ref{D:plane-matrix} to the first jump gives,
after testing it with $\phi(z)$, a fixed nonzero common constant times
\begin{equation}\label{D:jump}
 \left[\mu_n(q)(M^2+|r|^2)^2
 c_T^{(\bar n)}(M,r)\,\phi(-M^2)\right]_{j_A}(M_o).
\end{equation}
The normalization is worth checking.  The plane insertion is
$b\mu_n(q)e^{iq\cdot Y}$, conjugated by $U_a(u)$, and the input
polarization is $T_{\bar n\bar n}=a_i^2c_T^{(\bar n)}$ with
$a_i=(M^2+|r|^2)/(2b)$.  The normal Fourier measure is
$\dd a_x=\dd z/(2b)$.  Thus the jump with respect to $\dd z$
contains $2b^2a_i^2=(M^2+|r|^2)^2/2$; no boost factor remains.
The opposite channel has exactly the same constant.  These
identities can first be tested on external wave packets, so
\eqref{D:jump} means bounded operator application to an internal
$M\,\dd M$ wave and testing on compact output masses.

\Needspace{6\baselineskip}
\begin{lemma}[The subtraction ambiguities]\label{D:subtractions}
There are continuous $L^2(M\,\dd M)$-valued regularizations
$\chi^n_{\rm reg}(r,M)$ of
\begin{equation}\label{D:chi}
 \chi^n(r,M)=(M^2+|r|^2)c_T^{(\bar n)}(M,r)
\end{equation}
such that, with $w=q+r$ and
\begin{equation}\label{D:AB}
A_A=[\mu_n(q)\chi^n_{\rm reg}(r)]_{j_A^{(n)}},\qquad
B_A=[\mu_{\bar n}(r)\chi^{\bar n}_{\rm reg}(q)]_{j_A^{(\bar n)}},
\end{equation}
one has
\begin{equation}\label{D:two-horizon}
 (M_o^2+|w|^2)A_z-(|w|^2-M_o^2)B_z+2M_ow_iB_i
 =P_q+P_r.
\end{equation}
The degree of each polynomial is uniformly finite.  The coefficients
of $P_q$, polynomial in $q$, may depend on $r$; those of $P_r$,
polynomial in $r$, may depend on $q$.
\end{lemma}

\begin{proof}
Choose a positive subtraction point $z_0$ and retain the full finite
covariant tensor family.  Use one integer $d$ larger than all its
polynomial growth orders, perform the Cauchy subtraction on that
family, and then contract with $nn,\bar n\bar n$ and the output
frame.  The subtracted Cauchy kernel is
\[
 \frac{(z-z_0)^d}{(z'-z_0)^d(z'-z)}.
\]
Using \eqref{D:jump}, its value at $z=|r|^2$ is proportional to
\begin{equation}\label{D:cauchy-wave}
 \left[\mu_n(q)\chi^n(r,M)
 \left(\frac{|r|^2-z_0}{-M^2-z_0}\right)^d\right]_{j_A}.
\end{equation}
Subtracting this Cauchy transform leaves zero jump.  Its boundary
Fourier transform is supported in both closed half-lines, hence
at zero; the remainder is a polynomial in $z$.

Its dependence on the transverse parameter must also be controlled.
A null rotation with parameter $\beta$, fixing $n$, changes
\[
 r\longmapsto r+b\beta,\quad
 a\longmapsto a+\beta\cdot(q+r)+\tfrac12b|\beta|^2,\quad
 a_x\longmapsto a_x+\beta\cdot q.
\]
It fixes $b,M_o,q,X,z$ and transports the output adapted frame.
The remaining tensor contractions transform polynomially in
$\beta$, by a matrix independent of $z$.  This matrix commutes
with the common-order Cauchy subtraction.  Covariance of the split
commutator therefore makes the
coefficients of its subtraction polynomial polynomial in $r$ at
fixed $q$.  Evaluation at $z=|r|^2$ preserves this property.  The
opposite calculation gives a polynomial in $q$ at fixed $r$.
The degrees are uniform because the distribution order in a rest
frame and the degrees of the finite tensor representation are
finite.  The opposite calculation may hold its own boost coordinate
$a$ fixed; the evaluated contracted matrix has net boost weight zero,
so it can then be written at the same $b$ as the first calculation.

For completeness the wave regularization can be made explicitly.
Set $R_2=|r|^2$, $D=M^2+R_2$.  The change of spatial frame gives
\begin{equation}\label{D:polarization}
 c_T^{(\bar n)}=l/3+d_Ad_Bs_{AB},\qquad
 d_z=\frac{R_2-M^2}{D},\qquad d_i=-\frac{2Mr_i}{D}.
\end{equation}
Write $s_{ij}^{\rm TF}=s_{ij}+\delta_{ij}s_z/2$.  Expansion of
\eqref{D:polarization} gives
\begin{equation}\label{D:chi-expansion}
 \chi^n=\frac D3l+
 \left(R_2-5M^2+\frac{6M^4}{D}\right)s_z
 +\left(-4Mr_i+\frac{8M^3r_i}{D}\right)s_{zi}
 +\frac{4M^2r_ir_j}{D}s_{ij}^{\rm TF}.
\end{equation}
In particular the radial coefficient follows from
$s_{11}+s_{22}=-s_z$:
$((R_2-M^2)^2-2M^2R_2)/D=R_2-5M^2+6M^4/D$.
The vector coefficient is $2Dd_zd_i$.

Choose a smooth mass cutoff $\zeta$ equal to zero below $1$ and
one above $2$, and subtract $\zeta(M)$ times
\[
 \frac{M^2+R_2}{3}l+(M^2-5R_2)s_z
 +(4Mr_i-8r_iR_2/M)s_{zi}+4r_ir_js_{ij}^{\rm TF}.
\]
The remaining large-mass terms are $O(M^{-2})$ or better on
compact $r$ sets, and hence belong to $L^2(M\,\dd M)$.
At small mass, $|d|=1$ in \eqref{D:polarization}, so the norm is
bounded by a constant times $M^2+|r|^2$.  Dominated convergence
now proves continuity also at $r=0$; the same formulas give
polynomial growth in $r$.

Replacing \eqref{D:cauchy-wave} by this regularization changes
only an $L^2$-valued polynomial in $r$.  Indeed
\[
 (R_2-z_0)^d-(-M^2-z_0)^d
\]
is divisible by $D$, canceling the sole denominator in
\eqref{D:chi-expansion}.  Both regulated expressions are $L^2$
waves; polynomial interpolation then shows that every coefficient
of their polynomial difference is an $L^2$ wave.  Choose the same
mass cutoff and subtraction prescription in the other horizon.

Finally the output frame rotation is
\[
 j_z^{(n)}=\frac{|w|^2-M_o^2}{M_o^2+|w|^2}j_z^{(\bar n)}
 -\frac{2M_ow_i}{M_o^2+|w|^2}j_i^{(\bar n)}.
\]
The common gap value from Lemma~\ref{D:causal-continuation} and the
common jump normalization therefore give \eqref{D:two-horizon}
off $q=0$ and $r=0$.

There is no extra distribution supported on those axes.  The left
side is continuous in $q,r$, with values on compact mass tests.
Take polynomial interpolation projections $\mathcal P_q,
\mathcal P_r$ of degrees larger than the two degree bounds, using
nodes away from zero.  On the off-axis patch,
$(1-\mathcal P_q)(1-\mathcal P_r)$ annihilates the left side.
By continuity it does so everywhere.  Decomposing with these two
projections extends its separated-polynomial form to all $q,r$.
\end{proof}

Fix the transverse Fourier convention
\[
 \widehat g(r)=(2\pi)^{-2}\int e^{-ir\cdot z'}g(z')\,\dd^2z',
 \qquad g(z')=\int e^{ir\cdot z'}\widehat g(r)\,\dd^2r.
\]
If both position tests vanish near the origin, every term $P_q+P_r$
in \eqref{D:two-horizon} is annihilated.  For a real radial
$g\in C_c^\infty(\R^2\setminus\{0\})$, put
\begin{equation}\label{D:Phi}
 \Phi(g)=\int\widehat g(r)\chi^n(r,M)\,\dd^2r,
 \qquad \mu(g)=\int g(z')\mu(\dd z').
\end{equation}
The subtraction makes no difference in this formula.

Indeed, integrating a polynomial in $r$ against $\widehat g(r)$
gives a finite linear combination of derivatives of $g$ at zero,
all of which vanish.  With one such test in each transverse variable,
this removes $P_q$ and $P_r$ separately, even though their remaining
coefficients can depend on the other variable.  The regularized
formula supplies the meaning of this operation on the original wave.
To perform the horizon exchange explicitly, let $\mathsf R$ be the spatial
half-turn about $e_2$. It sends $n$ to $\bar n$, $\bar n$ to $n$, and acts
on transverse coordinates by $\mathsf S=\operatorname{diag}(-1,1)$.
The adapted tetrads obey
\[
 \mathsf R e_z^{(n)}(p)=e_z^{(\bar n)}(\mathsf Rp),\qquad
 \mathsf R e_i^{(n)}(p)=\mathsf S_{ij}e_j^{(\bar n)}(\mathsf Rp).
\]
Thus the longitudinal label has no extra sign. In the covariantly
identified output frames, with the same regularization on both horizons,
\[
 A_z(q,r)=B_z(\mathsf Sr,\mathsf Sq),\qquad
 A_i(q,r)=\mathsf S_{ij}B_j(\mathsf Sr,\mathsf Sq).
\]
Write $\langle F\rangle_g=\int\widehat g(q)\widehat g(r)F(q,r)
\,\dd^2q\,\dd^2r$. Radiality makes the test invariant under
$(q,r)\mapsto(\mathsf Sr,\mathsf Sq)$, including the weight
$|q+r|^2$. Consequently
\[
 \langle A_z\rangle_g=\langle B_z\rangle_g=:C_g,\qquad
 \langle |q+r|^2A_z\rangle_g=\langle |q+r|^2B_z\rangle_g,
\]
and \eqref{D:two-horizon} reduces to
\[
 2M_o^2C_g+2M_o\sum_i\langle(q_i+r_i)B_i\rangle_g=0.
\]
Our Fourier convention gives
$q_i\widehat g(q)=-i\widehat{\partial_i g}(q)$. Each of the two
summands differentiates one test, with the same sign. On returning
their vector contraction to the $n$ frame, the transverse matrix
$\mathsf S$ occurs twice and cancels. Division by $2M_o$ gives
\begin{equation}\label{D:radial-identity}
 M_o[\mu(g)\Phi(g)]_{j_z}
 =i\sum_i
 [\mu(\partial_i g)\Phi(g)+\mu(g)\Phi(\partial_i g)]_{j_i}.
\end{equation}
Reversing the Fourier orientation would change the common phase,
which will play no role; the relative coefficients are fixed.
In particular, \eqref{D:chi-expansion} shows that the radial and
vector kernels surviving away from the position origin are
$6M^4/(|r|^2+M^2)$ and $8M^3r_i/(|r|^2+M^2)$.
The two-dimensional massive Green transform gives respectively
$M^4K_0(M|z'|)$ and a fixed imaginary multiple of
$M^4\widehat z'_iK_1(M|z'|)$.  The scalar and transverse trace-free
terms do not contribute to a radial test, nor to its first angular
harmonic derivative, as appropriate.

\subsection{Exact Hankel leakage}

We now isolate the only obstruction to improving the order-two
comparison on $l$ to order zero.  A test with zero first radial moment
has no leakage beyond its Hankel support.  It follows that every
possible exterior tail depends on that single moment; dilation
covariance will determine its exact shape.

Let $\mathcal H_\nu$ be the unitary Hankel transform
\[
 (\mathcal H_\nu F)(M)=\int_0^\infty J_\nu(Mv)F(v)v\,\dd v.
\]
It diagonalizes $h_\nu$ as multiplication by $v^2$.
Put $\lambda=\sqrt\kappa$.  On its positive spectral part,
dilatation covariance permits a homogeneous realization
$L^2(\R_+,\lambda\,\dd\lambda;\mathcal K)$ with constant
multiplicity space $\mathcal K$.  Its zero spectral subspace is
kept separately.  Write $\mathcal U$ for this spectral
realization, and $f c$ for the input radial wave with label $c$.

\Needspace{6\baselineskip}
\begin{lemma}[Exact tail]\label{D:exact-tail}
For $c=l,j_z,s_z$ there is $d_c\in\mathcal K$ such that, if
$F\in C_c^\infty(0,R)$, then
\begin{equation}\label{D:tail}
 \mathcal U\bigl(c\mathcal H_0F\bigr)(\lambda)
 =\left(\int_0^\infty vF(v)\,\dd v\right)\lambda^{-2}d_c
 \quad\text{for almost every }\lambda>R.
\end{equation}
The label map is linear, and
\begin{equation}\label{D:tail-relations}
 d_{j_z}=-d_l,\qquad d_{s_z}=-d_l/3.
\end{equation}
\end{lemma}

\begin{proof}
The Bessel recurrence $J_2(x)=2J_1(x)/x-J_0(x)$ and
$\int_0^v rJ_0(Mr)\,\dd r=vJ_1(Mv)/M$ give
\begin{equation}\label{D:hankel-transfer}
 (\mathcal H_2\mathcal H_0F)(v)
 =\frac2{v^2}\int_0^v rF(r)\,\dd r-F(v).
\end{equation}
The identity first follows on smooth compact tests by the unitary
Hankel inversion and then in $L^2(v\,\dd v)$.  If $F$ has zero
first radial moment and support below $R$, the right side also
has support below $R$.  The order-two comparison in
\eqref{D:inherited-comparison} therefore puts
$c\mathcal H_0F$ in the $\lambda\leq R$ subspace.

Consequently the restriction of the transfer to any interval
strictly above $R$ factors through the single functional
$F\mapsto\int vF(v)\,\dd v$.  Choose one annular test with this
moment equal to one.  Overlapping support bounds identify the
resulting vector-valued coefficient at every $\lambda>0$.
In these coordinates internal dilatations act as
$f(M)\mapsto e^sf(e^sM)$ on the input and
$\psi(\lambda)\mapsto e^{-s}\psi(e^{-s}\lambda)$ on the output.
Since the moment of the corresponding Hankel-side test scales
by $e^s$, its coefficient must be $\lambda^{-2}d_c$.
A reference test and its square-integrable tail ensure
$d_c\in\mathcal K$.  This proves an exact tail, not only its
asymptotic order.

The order-zero comparisons on $c_V$ and $c_T$ say that
$c_V\mathcal H_0F$ and $c_T\mathcal H_0F$ have no spectrum above
the support radius of $F$.  Linearity of the tail gives
$d_l+d_{j_z}=0$ and $d_l/3+d_{s_z}=0$.
\end{proof}

\subsection{The logarithmic pairing}

The tail relations alone allow a nonzero vector $d_l$.  We will insert
a radial test covering many scales into \eqref{D:radial-identity}.
Its two scalar-channel tails produce a term proportional to
$\log L\,\|d_l\|^2$, whereas the first angular harmonics give smaller
terms.  The estimates below keep this leading coefficient and the
errors separate.

Choose $0\leq\eta_L\leq1$, smooth and supported in
$(0,\log L)$, equal to one outside fixed-width neighborhoods
of the endpoints, with uniformly bounded derivatives.  Set
\begin{equation}\label{D:plateau}
 g_L(z')=|z'|^2\eta_L(\log|z'|).
\end{equation}
We retain the nonzero normalization constant $c_*$ throughout
the estimates; it is determined by the radial Fourier kernel.

\Needspace{6\baselineskip}
\begin{lemma}[Plateau estimates]\label{D:plateau-estimates}
There are scalar functions $f_L,u_L$ such that
\begin{equation}\label{D:plateau-waves}
 \Phi(g_L)=s_zf_L(M),\qquad
 \Phi(\partial_i g_L)=M s_{zi}u_L(M),
\end{equation}
and, writing $\mathscr D=M\partial_M$,
\begin{align}
 &\norm{f_L-c_*\mathbf1_{(1/L,1)}}_{L^2(\dd M/M)}=O(1),
 &&c_*\ne0,\label{D:plateau-main}\\
 &\norm{\mathscr D^a f_L}_{L^2(\dd M/M)}+
 \norm{\mathscr D^a u_L}_{L^2(\dd M/M)}=O(1)
 &&(a=1,2),\label{D:plateau-derivatives}\\
 &\norm{\kappa\Phi(\partial_i g_L)}=O(1),\qquad
 \norm{f_L/M}_{L^2(M\,\dd M)}=O(\sqrt{\log L}).
 \label{D:plateau-norms}
\end{align}
For $c=s_z,j_z$, in the positive $\lambda$ representation,
\begin{equation}\label{D:plateau-tail}
 \lambda\mathcal U(c f_L)
 =c_*\mathbf1_{(1,L)}(\lambda)\lambda^{-1}d_c+E_{c,L},
 \qquad \norm{E_{c,L}}_{L^2(\lambda\,\dd\lambda;\mathcal K)}=O(1).
\end{equation}
\end{lemma}

\begin{proof}
The kernels computed after \eqref{D:radial-identity}, with
$t=M|z'|$, give explicitly in our transverse Fourier convention
\begin{align}
 f_L(M)&=6\int_0^\infty t^3K_0(t)
     \eta_L(\log t-\log M)\,\dd t,\label{D:f-convolution}\\
 u_L(M)&=-4i\int_0^\infty t^2K_1(t)
     (2\eta_L+\eta_L')(\log t-\log M)\,\dd t.
 \label{D:u-convolution}
\end{align}
Changing the Fourier orientation changes the fixed phase in the
second formula only.  The first has
\[
 c_*=6\int_0^\infty t^3K_0(t)\,\dd t=24
\]
with the convention fixed above.  In particular it cannot vanish.
On the logarithmic variable these are convolutions with integrable
kernels having finite first moment.  At $t=0$ their densities
decay exponentially in $|\log t|$, and at infinity they decay
faster; the relevant logarithmic derivatives have the same
integrability.  Convolution of the interval plateau differs from
its kernel integral times that plateau by a uniformly bounded
$L^2$ error: the difference is supported near the two edges up
to integrable tails.  Positive derivatives act on these two
fixed-width edges.  This proves
\eqref{D:plateau-main}--\eqref{D:plateau-derivatives}.
The functions themselves and their derivatives have integrable
tails at zero and infinity for each fixed $L$.

Direct differentiation gives
\[
 h_1(Mu)=-M^{-1}(\mathscr D^2+2\mathscr D)u,
 \qquad
 \norm{h_1(Mu)}^2=
 \int_0^\infty|\mathscr D^2u+2\mathscr Du|^2\frac{\dd M}{M}.
\]
The order-one comparison proves the first estimate in
\eqref{D:plateau-norms}; the second is
\eqref{D:plateau-main}.

To prove \eqref{D:plateau-tail} without assigning a divergent
form norm to an individual nonzero-moment test, fix a smooth
annular $F$ with $\int vF(v)\,\dd v=1$ and put $G=\mathcal H_0F$.
Replace $f_L$ by
\[
 c_*\mathcal H_0\bigl(F-L^{-2}F(\cdot/L)\bigr)(M)
 =c_*\bigl(G(M)-G(LM)\bigr).
\]
Since $G(0)=1$ and $G$ decreases rapidly at infinity, this
difference has the same plateau and the same uniform edge
estimates as $f_L$.  Its error $e_L$ satisfies
\[
 \norm{\sqrt{h_2}e_L}^2
 =\int_0^\infty
 (|\mathscr De_L|^2+4|e_L|^2)\frac{\dd M}{M}=O(1).
\]
The comparison \eqref{D:inherited-comparison} gives an $O(1)$
error after multiplication by $\lambda$.  Above a fixed radius,
the $G$ term has the exact tail of Lemma~\ref{D:exact-tail}.
The norm of $G(LM)$ in $L^2(M\,\dd M)$ is $L^{-1}\norm G$;
hence its $\lambda$-weighted norm below $CL$ is $O(1)$ for
any fixed $C$.  Above both Hankel support radii the two exact
tails cancel, because their moments are both one.  The remaining
fixed lower interval and fixed-ratio upper interval contribute
$O(1)$ to the weighted norm.  This proves
\eqref{D:plateau-tail}.  The $\lambda=0$ spectral subspace
contributes nothing to that estimate.

Using the difference of the two dilated tests is essential here:
their individual $\lambda^{-2}$ tails need not have finite
$\lambda$-weighted norm, but the exact cancellation above the larger
support radius makes their difference admissible.  No estimate has
been assigned to either divergent weighted tail separately.
\end{proof}

\Needspace{6\baselineskip}
\begin{proposition}[Order-zero scalar comparison]
\label{D:critical-comparison}
The quadratic moment $\kappa$ of the physical bounded-profile
generators satisfies
\begin{equation}\label{D:order-zero}
 \mathbf1_{(s,\infty)}(\kappa)
 \mathbf1_{[0,s]}(h_0)\big|_l=0\qquad(s\geq0).
\end{equation}
\end{proposition}

\begin{proof}
On the positive $\lambda$ spectrum define
\[
 A_L(\lambda)=\lambda^{-2}\mu(g_L),\qquad
 B_{i,L}(\lambda)=\lambda^{-1}\mu(\partial_i g_L).
\]
The moments in \eqref{D:measure-input} imply
$0\leq A_L\leq\id$ and $\norm{B_{i,L}}\leq C$ uniformly:
$|\partial_i g_L(z')|\leq C|z'|$, and positive-measure
Cauchy--Schwarz bounds its integral by $C\sqrt\kappa$.
Scale covariance identifies the law of $z'/\lambda$ in each
fiber with one fixed positive measure on $\mathcal K$.  Its
mean is zero and its trace-second moment is $\id$.  Hence,
on every fixed multiplicity vector,
\begin{equation}\label{D:bulk}
 A_L(\lambda)\longrightarrow\id,\qquad
 B_{i,L}(\lambda)\longrightarrow0
 \quad\text{if }\lambda\longrightarrow\infty,
 \quad L/\lambda\longrightarrow\infty.
\end{equation}
For the first limit use the integrable weight $|z'/\lambda|^2$;
for the second use $|z'/\lambda|$ and the zero first moment.
The limiting derivative integrand is $2z'_i/\lambda$.
Bounded convergence for the positive measure proves strong
convergence; no uniform rate is required.

These are statements about the fixed multiplicity vectors that enter
the tail formula.  Positivity and the second moment control the
truncated integrands; projection-valuedness of $\mu$ is unnecessary.
On $\ker\kappa$ the second moment is zero, so $\mu$ is concentrated at
$z'=0$ there.  The profiles $g_L$ and $\partial_i g_L$ vanish near that
point, and hence this spectral subspace makes no contribution to the
pairing.

Pair \eqref{D:radial-identity} with $f_L(M_o)/M_o$ and use
\eqref{D:plateau-waves}.  This is allowed although the identity
was first tested on compact mass intervals.  Cut off this test
smoothly at zero and infinity.  The approximants converge both
in its $L^2(M_o\,\dd M_o)$ norm and after multiplication by
$M_o$.  For fixed $L$, the profile operators are bounded, and
\eqref{D:plateau-norms} controls the derivative wave, so every
pairing has a limit.

The left side becomes
\begin{align*}
 \ip{j_zf_L}{\mu(g_L)s_zf_L}
 &=\ip{\lambda\mathcal U(j_zf_L)}
 {A_L\lambda\mathcal U(s_zf_L)}\\
 &=|c_*|^2\int_1^L
 \ip{d_{j_z}}{A_L(\lambda)d_{s_z}}\frac{\dd\lambda}{\lambda}
 +O(\sqrt{\log L}).
\end{align*}
The exact normalization $|c_*|^2$ is common to the bra and ket.
Equation~\eqref{D:bulk} and boundedness imply, by averaging on
the logarithmic interval, that this is
\begin{equation}\label{D:leading-pairing}
 |c_*|^2\log L\,\ip{d_{j_z}}{d_{s_z}}+o(\log L).
\end{equation}
For example, discard two logarithmic edge intervals of fixed
width, apply the strong bulk convergence on the remaining
interval, and then let the discarded width grow.  This proves
the asserted Cesaro limit directly.

For the first term on the right of \eqref{D:radial-identity},
\eqref{D:bulk}, \eqref{D:plateau-tail} and uniform boundedness give
\[
 \norm{\mu(\partial_i g_L)s_zf_L}
 =\norm{B_{i,L}\lambda\mathcal U(s_zf_L)}
 =o(\sqrt{\log L}).
\]
The output test $j_if_L/M_o$ has norm $O(\sqrt{\log L})$.
Their pairing is therefore $o(\log L)$.  For the last term,
$0\leq\mu(g_L)\leq\kappa$ and commutativity give
\[
 \norm{\mu(g_L)\Phi(\partial_i g_L)}
 \leq\norm{\kappa\Phi(\partial_i g_L)}=O(1).
\]
Its pairing is $O(\sqrt{\log L})$.
Comparison with \eqref{D:leading-pairing} forces
\[
 0=\ip{d_{j_z}}{d_{s_z}}=\tfrac13\norm{d_l}^2,
\]
by \eqref{D:tail-relations}.  Thus $d_l=0$.
If either of the spin labels used above is zero, the same
conclusion already follows from \eqref{D:tail-relations}.

Lemma~\ref{D:exact-tail} now says that $l\mathcal H_0F$ has no
$\lambda$ spectrum above $R$ whenever $F$ is supported below
$R$.  Smooth compact tests in $(0,R)$ are dense in
$L^2((0,R),v\,\dd v)$.  Hankel unitarity and closedness of
spectral subspaces give \eqref{D:order-zero} with $s=R^2$;
at $s=0$ the regular $h_0$ copy has no zero eigenvector.
The result concerns the moment $\kappa$.  It does not identify
the full parabolic operator $h$ with that moment.
\end{proof}
